\textbf{Gene-deletion explanations in biological context.}
\textbf{a,b,} RNA UMAPs show reference labels and prediction errors (red) for all 320 PBMC queries (100 supports/class) and 7,769 development-cortex queries (ten supports/class), both draw 0. Paired views share coordinates fitted without labels to the support/query union (Methods).
\textbf{c,} PBMC support mean log1p(counts per 10,000), detection fraction and mean corrected DS-LOO query attribution. Independent scales distinguish expression, detection and classifier-margin units. The 19 genes combine each class's first-ranked gene and all marker checks in Supplementary Table~\ref{tab:s8}. Positive attribution supports the predicted class relative to its fixed runner-up; negative attribution opposes that margin.
\textbf{d,} All 20,365 exported PBMC cell--gene validation pairs, including 1,625 additional DS-LOO-selected pairs and retaining zero and negative effects; dashed line, identity. Exact deletion covers selected genes, not every gene.
\textbf{e,} DS-LOO and IG64 fidelity on the same original evaluated gene set: within-cell Spearman correlations summarized within and across classes. Labels give supports/draw. Six configurations, including two kidney representations, are not six independent cohorts.
\textbf{f,} Ranked deletion versus two pooled random-expressed orderings/cell reverses 45.0\% versus 0.3\% of original margins at 50 genes.
Explanations use 191/320 PBMC and 96/7,769 cortex queries correctly classified in both stored and float32 predictions. Summaries are descriptive, without confidence intervals. Zeroing an absent gene cannot assess absence information; model-deletion fidelity does not establish biological causality.
